\documentclass[10pt]{article}
\usepackage[T1]{fontenc}
\usepackage{lmodern,amsmath,amssymb,amsthm,microtype}
\usepackage[margin=.8in]{geometry}
\usepackage[colorlinks=true,urlcolor=blue]{hyperref}
\newtheorem{theorem}{Theorem}
\title{The Rational Generating Function of A308620}
\author{Qichao Wang\\{\small Hebei University of Technology}\\[.3em]Daoyu Dong\\{\small University of Electronic Science and Technology of China}}
\date{October 26, 2026}
\begin{document}\maketitle
\begin{abstract}
We derive the generating function conjectured for the steady-state
counts of the evolutionary spatial prisoner's dilemma on a path.
The proof starts from the established exact count, rather than a
fit to finitely many values. It proves the proposed order-ten
recurrence with its correct starting index: the recurrence applies
from index 12, and fails by one at index 11.
\end{abstract}
\section{The exact counting formula}
Write $a(n)$ for A308620~\cite{oeis}, $n\ge1$. The counts are for
path states modulo reflection. Burger, van der Merwe and van Vuuren
prove the following formula in Theorem 3, equation (4.2) of
\cite{paper}; it is also the entry's Maple evaluator:
\[
 a(n)=1+Q_n+\tfrac12(R_n+\bar R_n+S_n+\bar S_n),
\]
where, with empty sums zero,
\begin{align*}
 Q_n&=\sum_{i=1}^{\lfloor(n+1)/5\rfloor}\binom{n-3i}{2i-1},\\
 R_n&=\sum_{i=1}^{\lfloor(n+4)/5\rfloor}\binom{n-3i+2}{2i-2},
 &\bar R_n&=\sum_{i=1}^{\lfloor(n+4)/5\rfloor}
 \binom{\lfloor(n-5i+4)/2\rfloor+i-1}{i-1},\\
 S_n&=\sum_{i=1}^{\lfloor(n-2)/5\rfloor}\binom{n-3i-2}{2i},
 &\bar S_n&=\sum_{i=1}^{\lfloor(n-2)/5\rfloor}
 \binom{\lfloor(n-5i-2)/2\rfloor+i}{i}.
\end{align*}
This established theorem supplies the connection to the game:
nontrivial steady states have defection runs of length at least
two, internal cooperation runs of length at least three, and
endpoint cooperation runs of length at least two. The terms above
count their run-length compositions and reflection-fixed cases.
We prove the subsequent generating-function conjecture, not a new
claim of authorship for this state classification.

\section{Summation and recurrence}
Put $D=(1-x)^2-x^5$, $E=1-x^2-x^5$, and let each capital letter
also denote the generating function over $n\ge1$ of its subscripted
sequence. The elementary identity
$\sum_{j\ge0}\binom{j+h}{h}x^j=(1-x)^{-h-1}$ gives
\begin{align*}
 Q&=\sum_{i\ge1}\frac{x^{5i-1}}{(1-x)^{2i}}=\frac{x^4}{D},
 &R&=\frac{x(1-x)}{D},&S&=\frac{x^7}{(1-x)D},\\
 \bar R&=\frac{x(1+x)}{E},
 &\bar S&=\frac{x^7}{(1-x)E}.
\end{align*}
For the barred sums, pair the two powers $x^{2j}$ and $x^{2j+1}$
to obtain
$\sum_{j\ge0}\binom{\lfloor j/2\rfloor+h}{h}x^j
=(1+x)/(1-x^2)^{h+1}$.
All the sums are formal: every coefficient receives finitely many
contributions, so rearrangement requires no analytic convergence.

\begin{theorem}
The ordinary generating function is exactly
\[
 G(x)=\sum_{n\ge1}a(n)x^n
 =2x-\frac{x^2(-2+2x+x^2-2x^3+3x^5+x^7+x^9)}
 {(x^5+x^2-1)(x^5-x^2+2x-1)}.
\]
For every $n\ge12$,
\[
 a(n)=2a(n-1)-2a(n-3)+a(n-4)+2a(n-5)
       -2a(n-6)-a(n-10).
\]
At $n=11$ the right side is 41 whereas $a(11)=42$.
\end{theorem}
\begin{proof}
Substitute the five rational expressions above into
$G=x/(1-x)+Q+(R+\bar R+S+\bar S)/2$ and combine fractions.
An equivalent polynomial identity, which also certifies the algebra,
is
\begin{align*}
 &(1-2x+2x^3-x^4-2x^5+2x^6+x^{10})G(x)\\
 &\hspace{1cm}=2x-2x^2-2x^3+3x^4-4x^6+x^7-x^9+x^{11}.
\end{align*}
The denominator is $DE$. Coefficients above degree 11 give the
recurrence. The coefficient at degree 11 is one, explaining its
explicit exception there. This proves both official conjectures
with the required initial-index qualification.
\end{proof}
\paragraph{Validation.}
The verifier computes the original binomial formula through
$n=1000$, checks every recurrence instance and the entire rational
polynomial identity in that range, and compares all 40 published
terms. A separate binary-word enumeration through length 18 checks
the run conditions and counts reversal orbits directly, independently
of the binomial formula. Finite checks validate the implementation;
the formal identities prove the all-index claims. No OEIS submission
or literature-priority claim is made.
\begin{thebibliography}{2}
\bibitem{oeis} R. J. Mathar, OEIS A308620,
\url{https://oeis.org/A308620}, accessed October 26, 2026.
\bibitem{paper} A. P. Burger, M. van der Merwe and J. H. van Vuuren,
An asymptotic analysis of the evolutionary spatial prisoner's dilemma
on a path, Discrete Applied Mathematics 160 (2012), 2075--2088,
Theorem 3 and equation (4.2), author-hosted copy:
\url{https://www.vuuren.co.za/papers/GamesPaths.pdf}.
\end{thebibliography}\end{document}
